\sectioncentered*{Заключение}
\addcontentsline{toc}{section}{Заключение}

Для трёх классов сообщений спроектированы девять логических правил с приоритетами и девять классов фраз (по две шаблонные фразы в каждом). Правила опираются на класс сообщения и сущность \texttt{rrel\_entity}. Один и тот же вопрос об исполнителе, авторе или альбоме даёт разные ответы в зависимости от того, какие факты есть в базе знаний. Ответы видны в интерфейсе NIKA.

\begin{thebibliography}{9}
	\addcontentsline{toc}{section}{Список использованных источников}

	\bibitem{nika}
	NIKA: Intelligent Knowledge-driven Assistant [Электронный ресурс]. \mdash  Режим доступа: \url{https://github.com/ostis-apps/nika}. \mdash  Дата доступа: 01.10.2026.

	\bibitem{ostis}
	OSTIS: Open Semantic Technology for Intelligent Systems [Электронный ресурс]. \mdash  Режим доступа: \url{https://ostis.net/}. \mdash  Дата доступа: 01.10.2026.

	\bibitem{scs}
	SC-код. Руководство пользователя [Электронный ресурс]. \mdash  Режим доступа: \url{https://ostis-dev.github.io/sc-machine/}. \mdash  Дата доступа: 01.10.2026.
\end{thebibliography}
